\section{Which Occupational Comparisons Support the Pattern?}\label{sec:results}

\subsection{Public benchmark reconstruction}\label{subsec:public_benchmark}

We first reconstruct the public stock comparison before conditioning on broad
occupation families.  These are independent CPS analogues using equal-occupation
Rule-A beta quintiles on the fixed primary support, not an exact replication of
BCC: exhaustive BCC occupation membership is unavailable.  For all employed
workers, the annual 2022--2024 Q5--Q1 growth-factor difference is $-0.0346$
($[-0.1455,0.0763]$); the November 2022--June 2026 contrast is $-0.2312$
($[-0.4660,0.0036]$).  The top-two/bottom-three dashboard analogue is a
different grouping, with estimates $-0.0269$ ($[-0.0972,0.0435]$) and
$-0.1007$ ($[-0.2380,0.0367]$), respectively.  Its inference is computed
separately; no Q5--Q1 rejection or precision statement transfers to it.

No-control occupation long differences provide another distinct estimand.  With
baseline employment-stock weights and HC1 standard errors, the all-employed
Q5--Q1 coefficient is $-0.0318$ ($[-0.1447,0.0811]$) over 2022--2024 and
$-0.2466$ ($[-0.4771,-0.0160]$) from November 2022 to June 2026.  On the
468-occupation common support, the separately defined no-Webb young-relative
model gives a pooled coefficient of $-0.1146$ ($[-0.2057,-0.0236]$); allowing
family-specific young-relative calendar-month paths gives $-0.0124$
($[-0.1446,0.1197]$).  Family conditioning therefore changes the descriptive
estimand by absorbing cross-family evolution.  It does not identify a causal AI
effect, and intervals containing zero establish non-detection rather than
equivalence.

An independent annual ACS reconstruction gives a different view of sampling
precision.  Between 2022 and 2024, the equal-occupation BCC analogue has an
all-employed Q5--Q1 growth-factor difference of $-0.0859$
($[-0.1238,-0.0480]$ under ACS successive-difference replication) and a
full-time civilian wage-and-salary difference of $-0.0834$
($[-0.1272,-0.0396]$).  The corresponding published BCC comparisons are
$-0.022$ and $-0.019$.  Because exhaustive BCC occupation membership is not
public, the gap cannot be assigned separately to samples and grouping rules.

The annual young-relative ACS extension also remains sensitive to broad-family
conditioning and calendar construction.  On the fixed primary grouping and the
full 2017--2019, 2021, 2023--2024 calendar, the all-employed coefficient moves
from $-0.1089$ ($[-0.1310,-0.0869]$) pooled to $-0.0474$
($[-0.0885,-0.0063]$) with family-year effects.  On a nonreuse calendar that
drops 2018--2019, the corresponding estimates are $-0.0951$
($[-0.1204,-0.0698]$) and $-0.0220$ ($[-0.0705,0.0264]$).  These intervals
capture ACS person-sampling uncertainty under a year-block approximation.
Occupation- and family-shock intervals are wider and include zero for each of
the four reported family-year variants.  The ACS therefore supplies a larger
public-sample analogue, not an exact payroll replication or causal AI design.

Table~\ref{tab:occupation} begins with the rebuilt pooled result.  The Q5--Q1 coefficient is $-0.1321$, with standard error 0.0452 clustered by occupation.  Its 95-percent interval is $[-0.2206,-0.0437]$.  Giving every SOC two-digit family its own young-relative calendar-month path yields $-0.0217$ ($[-0.1607,0.1173]$).  On common draws, the conditioned-minus-pooled movement is 0.1104 ($[0.0107,0.2102]$), with paired $\mathrm{MDE}_{80}=0.1454$.  Thus the conditional estimands are distinguishable while the residual gradient is imprecise.

\begin{table}[!htbp]
\centering
\caption{Pooled and broad-family exposure comparisons}\label{tab:occupation}
\begin{tabular}{lrrrr}
\toprule
Specification & Coefficient & 95\% interval & Paired movement & $\mathrm{MDE}_{80}$ \\
\midrule
Pooled Q5--Q1 & $-0.1321$ & $[-0.2206,-0.0437]$ & --- & 0.1266 \\
SOC2 $\times$ young $\times$ month & $-0.0217$ & $[-0.1607,0.1173]$ & $0.1104$ & 0.1998 \\
\bottomrule
\end{tabular}
\tabnote{The weighted-stock criterion is the grouped-binomial conditional likelihood for the two age-specific employment-stock counts.  Coefficients parameterize log ratios of conditional means; no observed log ratio is constructed, and one-sided zero-stock cells remain in the criterion.  The second row changes the conditioning estimand.  Intervals use 9,999 occupation wild-score draws.  The paired movement is the second-row coefficient minus the first on common draws; its interval is $[0.0107,0.2102]$ and its $\mathrm{MDE}_{80}$ is 0.1454.  MDEs are normal-theory precision diagnostics, not rejection thresholds.}
\end{table}


The ratio $0.1104/0.1321$ is not a causal ``share explained.''  Broad families differ in AI-relevant tasks, computerization, education, work location, and pandemic recovery; family paths absorb all such common evolution.  The result says that the pooled estimate relies mainly on comparisons across families, not what would have happened without AI.

The full Q2--Q5 profile reaches the same conclusion.  Pooled coefficients are jointly different from zero but do not establish a monotone dose response.  After family-month conditioning, the four coefficients are jointly indistinguishable from zero.  Simultaneous intervals and paired movements appear in the online appendix.

\subsection{Direct support and restrictions}

Only SOC27, 29, 31, and 41 contain Q1 and Q5 occupations.  Elsewhere the regression links partial segments---for example Q1--Q3 in one family and Q4--Q5 in another---through common quintile coefficients.  That is valid regression variation but not repeated side-by-side tail comparisons.

\begin{table}[!htbp]
\centering
\caption{Exposure-quintile support within broad occupation families}\label{tab:family_quintile_support}
\scriptsize
\setlength{\tabcolsep}{2.4pt}
\resizebox{\textwidth}{!}{%
\begin{tabular}{r l ccccc c}
\toprule
SOC2 & Occupation family & Q1 & Q2 & Q3 & Q4 & Q5 & Both tails \\
\midrule
11 & Management & 0/0.00 & 0/0.00 & 15/4.85 & 12/8.45 & 2/0.25 &  \\
13 & Business and financial & 0/0.00 & 0/0.00 & 1/0.74 & 16/2.42 & 11/3.06 &  \\
15 & Computer and mathematical & 0/0.00 & 0/0.00 & 0/0.00 & 1/0.03 & 14/4.29 &  \\
17 & Architecture and engineering & 0/0.00 & 0/0.00 & 3/0.12 & 9/1.49 & 7/0.67 &  \\
19 & Life, physical, and social science & 0/0.00 & 0/0.00 & 12/0.86 & 6/0.34 & 3/0.05 &  \\
21 & Community and social service & 1/0.03 & 5/0.59 & 6/0.69 & 0/0.00 & 0/0.00 &  \\
23 & Legal & 0/0.00 & 0/0.00 & 1/0.91 & 2/0.35 & 1/0.08 &  \\
25 & Education and library & 0/0.00 & 3/3.52 & 2/0.85 & 0/0.00 & 2/0.16 &  \\
27 & Arts, media, and sports & 2/0.03 & 5/0.27 & 3/0.27 & 6/0.57 & 9/0.65 & Yes \\
29 & Healthcare practitioners & 3/0.20 & 18/1.94 & 11/3.80 & 3/0.21 & 2/0.23 & Yes \\
31 & Healthcare support & 4/1.58 & 6/1.45 & 3/0.61 & 0/0.00 & 1/0.03 & Yes \\
33 & Protective service & 3/0.30 & 8/1.05 & 6/0.92 & 0/0.00 & 0/0.00 &  \\
35 & Food preparation and serving & 5/1.14 & 4/1.76 & 0/0.00 & 1/0.37 & 0/0.00 &  \\
37 & Building and grounds maintenance & 3/1.96 & 2/0.34 & 1/0.20 & 0/0.00 & 0/0.00 &  \\
39 & Personal care and service & 7/1.72 & 8/0.56 & 2/0.23 & 0/0.00 & 0/0.00 &  \\
41 & Sales and related & 1/0.03 & 2/1.53 & 3/2.75 & 4/3.48 & 7/2.28 & Yes \\
43 & Office and administrative support & 0/0.00 & 5/0.80 & 3/0.04 & 7/2.60 & 32/7.77 &  \\
45 & Farming, fishing, and forestry & 3/0.09 & 2/0.05 & 0/0.00 & 0/0.00 & 0/0.00 &  \\
47 & Construction and extraction & 27/5.49 & 1/0.02 & 2/0.62 & 0/0.00 & 0/0.00 &  \\
49 & Installation, maintenance, and repair & 21/2.65 & 8/0.56 & 2/0.35 & 0/0.00 & 0/0.00 &  \\
51 & Production & 39/2.13 & 9/0.72 & 4/1.28 & 0/0.00 & 0/0.00 &  \\
53 & Transportation and material moving & 10/2.69 & 13/4.77 & 2/0.06 & 0/0.00 & 0/0.00 &  \\
\bottomrule
\end{tabular}}
\tabnote{Each Q cell is occupation count/national preperiod employment-stock share in percent.  The producer forms each occupation's stock by summing monthly CPS-weighted young and older employment over the 71 observed months from January 2017 through November 2022.  Because every occupation contributes the same months, these shares also equal shares based on average-month stock.  They sum to 100 percent across the 110 cells; zero cells are shown rather than suppressed.  The matrix uses the fixed 468-occupation support.  Only four families contain both Q1 and Q5.  The pooled regression connects other families through intermediate quintiles and common coefficients, but they do not supply direct within-family tail comparisons.  Family totals, beta ranges, and information shares are in Online Appendix Table C.1.}
\end{table}


A direct-tail benchmark uses only Q1 and Q5 occupations in the four spanning families.  Its 29 occupations account for 5.03 percent of preperiod stock.  The family-month estimate is 0.1494 ($[-0.1694,0.4682]$), so the CPS cannot precisely resolve a direct tail contrast within broad families.  This row changes population and support and cannot replace the primary estimand.

A continuous companion instead subtracts each family's preperiod-weighted beta mean and imposes a common within-family slope.  Its family-month coefficient is $-0.0025$ per 0.1084 raw beta units ($[-0.0240,0.0191]$).  It uses more residual variation but trades the tail contrast for linearity and a common slope.  Together, the specifications expose the support--functional-form tradeoff.

The pre-2020 occupational bridge adds a separate allocation uncertainty.  A
jointly feasible odds-tilt grid reallocates the 14.86 percent of supported
early-period stock routed through 31 split source occupations while preserving
source--age--month mass.  With exposure groups rebuilt after each allocation,
the pooled coefficient ranges from $-0.1446$ to $-0.1163$, and all six paired
intervals relative to the no-tilt allocation contain zero.  The family-month
estimand is less stable: at the $K=0.25$ endpoint it moves from $-0.0217$ to
$0.0439$, a paired change of $0.0656$ ($[0.0031,0.1280]$), and five of six
nonbaseline paired intervals exclude zero.  A wider 25-point feasible grid
places the pooled coefficient between $-0.1599$ and $-0.0702$.  These are
post-outcome mapping sensitivities over finite declared grids, not estimates of
true age-specific routes, sharp bounds, or identified sets.

\subsection{Influence and observed paths}

The online appendix reports nuisance-adjusted information, effective information counts, and family and occupation deletions.  Fast-food and counter workers have substantial signed influence because low-exposure service employment recovered differently after the pandemic, but other occupations push oppositely.  Outcome-selected deletion is not a design correction; symmetric trimming and economically defined service exclusions remain diagnostics.

Figure~\ref{fig:paths} separates young and older stocks for both tails.  Exact
equal-month accounting gives a Q5--Q1 young-stock change of $-0.1098$ log point
and an older-stock change of $0.0483$, so their young-relative difference is
$-0.1581$.  The numerator supplies the larger movement, while the denominator
reinforces it.  This identity is not the grouped-binomial coefficient
$-0.1321$, and neither component is an individual probability or a causal
share.  A symmetric family accounting assigns $-0.1293$ to within-family
ratio changes and $-0.0288$ to changing older-family weights; the appendix
reports the zero-denominator audit.  The low-exposure reference trajectory is
part of every comparison.

\begin{figure}[!htbp]
\centering
\includegraphics[width=0.94\textwidth]{r3_figure1_q1_q5_paths.pdf}
\caption{Young and older employment stocks in the exposure tails}\label{fig:paths}
\tabnote{Monthly CPS-weighted stocks use rebuilt beta labels and are indexed to each series' 2019 mean.  The ratio panel divides young by older stock before indexing.  The vertical line marks January 2023; December 2022 is the transition month, and October 2025 is absent rather than interpolated.  Paths are descriptive.}
\end{figure}
